\chapter{Revisão da Literatura}
\label{cap:revisao}

\section{Modelos de Fatores: dos Estáticos aos Latentes Dinâmicos}
\label{sec:modelos-fatores}

\begin{itemize}
    \item \textbf{Estrutura geral do modelo de fatores:} retorno de um ativo como combinação linear de exposições a fatores de risco comuns mais componente idiossincrático; equação que será retomada com notação completa no Capítulo~\ref{cap:metodologia}:
\end{itemize}

\begin{equation}
    r_{i,t} = \alpha_i + \sum_{k=1}^{K} \beta_{i,k}\, f_{k,t} + \varepsilon_{i,t}
    \label{eq:modelo-fatorial}
\end{equation}

\begin{itemize}
    \item \textbf{CAPM e APT (primeira geração):} CAPM \parencite{sharpe1964} --- fator único (excesso de retorno do portfólio de mercado), betas estimados por regressão de série temporal; APT \parencite{ross1976} --- múltiplos fatores sob ausência de arbitragem, sem derivação de equilíbrio. Betas constantes; fatores definidos ex-ante pelo pesquisador. Ganho conceitual: estrutura formal que relaciona retornos a fontes de risco sistemático.
    \item \textbf{Modelos de características observáveis (segunda geração):} Fama-French três fatores \parencite{fama1993} --- mercado, SMB (tamanho) e HML (valor); Fama-French cinco fatores \parencite{fama2015} --- acrescenta RMW (rentabilidade) e CMA (investimento); Carhart \parencite{carhart1997} --- acrescenta WML (momento). Ganho conceitual: fatores passam a ter interpretação econômica como prêmios de risco associados a características observáveis; ainda lineares, estáticos e baseados em hipóteses ex-ante.
    \item \textbf{Proliferação e o \textit{factor zoo}:} centenas de fatores documentados na literatura sem critério de seleção sistemático \parencite{harvey2016} expõem a fragilidade da especificação ex-ante e motivam abordagens que aprendem os fatores dos dados.
    \item \textbf{IPCA --- betas variantes em características observáveis \parencite{kelly2019}:} exposições fatoriais $\beta_{i,t}$ como função linear das características dos ativos; variam no tempo quando as características variam; estimação por mínimos quadrados alternados. Ganho conceitual: betas deixam de ser constantes no tempo; relação características-betas ainda é linear.
    \item \textbf{Conditional Autoencoder --- não-linearidades aprendidas \parencite{gu2021}:} fatores latentes extraídos por autoencoder; exposições como função não-linear das características dos ativos, aprendidas conjuntamente; supera sistematicamente modelos lineares e redes neurais sem estrutura fatorial no \textit{cross-section} norte-americano. Ganho conceitual: especificação ex-ante dos fatores é eliminada; não-linearidades são capturadas. Limitação remanescente: produz estimativas pontuais de retorno esperado sem modelar a distribuição condicional completa.
    \item \textbf{Posicionamento do \textit{FactorVAE}:} o salto seguinte --- acrescentar representação probabilística dos fatores latentes --- é o que distingue \textcite{duan2022} do Conditional Autoencoder; a formalização completa é apresentada no Capítulo~\ref{cap:metodologia}.
\end{itemize}

% PLACEHOLDER: figura ilustrando a progressão histórica --- linha do tempo ou diagrama com
% os cinco estágios (CAPM/APT -> FF/Carhart -> IPCA -> Conditional Autoencoder -> FactorVAE)
% e o ganho conceitual em cada salto.

\section{Variational Autoencoders}
\label{sec:vae}

\begin{itemize}
    \item \textbf{Autoencoder determinístico vs.\ variacional:} autoencoder determinístico aprende uma função de reconstrução ponto a ponto --- dado $x$, produz um único ponto $z$ no espaço latente; VAE aprende uma \emph{distribuição} sobre o espaço latente --- dado $x$, o codificador produz parâmetros $(\mu, \sigma)$ de uma gaussiana $q(z|x)$, e $z$ é amostrado dessa distribuição.
    \item \textbf{Por que a representação probabilística é útil em dados ruidosos:} em dados financeiros com baixa razão sinal-ruído, a incerteza sobre os fatores latentes é informação relevante; tratar os fatores como variáveis aleatórias permite capturar e propagar essa incerteza; o modelo pode expressar que "os fatores são ambíguos hoje" em vez de forçar uma estimativa pontual possivelmente enganosa.
    \item \textbf{O problema da inferência:} a log-verossimilhança marginal $\log p(x) = \log \int p(x|z)\,p(z)\,\mathrm{d}z$ é intratável quando o decodificador $p(x|z)$ é parametrizado por rede neural não-linear --- a integral não tem forma fechada.
    \item \textbf{ELBO e inferência variacional:} a solução é introduzir uma distribuição aproximada $q(z|x)$ (codificador) e maximizar um limite inferior tratável da log-verossimilhança:
\end{itemize}

\begin{equation}
    \log p(x) \;\geq\; \underbrace{\mathbb{E}_{q(z|x)}\!\left[\log p(x|z)\right]}_{\text{qualidade de reconstrução}} \;-\; \underbrace{\mathrm{KL}\!\left(q(z|x) \,\|\, p(z)\right)}_{\text{regularização (divergência KL)}}
    \label{eq:elbo}
\end{equation}

\begin{itemize}
    \item \textbf{Interpretação dos dois termos:} o primeiro maximiza a probabilidade de reconstruir $x$ a partir de amostras de $z \sim q(z|x)$; o segundo penaliza o quanto a posterior aprendida $q(z|x)$ se afasta da prior $p(z)$, evitando colapso do espaço latente e forçando representação compacta.
    \item \textbf{Divergência KL:} quando $q(z|x) = \mathcal{N}(\mu_\phi, \sigma_\phi^2 I)$ e $p(z) = \mathcal{N}(0, I)$, o KL tem forma fechada analítica --- não requer amostragem e é computacionalmente barato.
    \item \textbf{Truque de reparametrização \parencite{kingma2014}:} para que o gradiente flua através da amostragem de $z \sim q(z|x)$, reescreve-se $z = \mu_\phi + \sigma_\phi \odot \epsilon$, com $\epsilon \sim \mathcal{N}(0,I)$; o gradiente flui por $\mu_\phi$ e $\sigma_\phi$, não pela operação de amostragem.
    \item \textbf{Prior aprendida:} quando a prior $p(z)$ é parametrizada por uma rede separada (como no \textit{FactorVAE}), o KL compara duas gaussianas com parâmetros distintos e os gradientes fluem de volta à rede que parametriza a prior --- esse ponto é diretamente relevante para o esquema professor-aluno do modelo.
    \item \textbf{Nível de apresentação:} ênfase em por que a representação probabilística importa e em como o ELBO torna o treinamento tratável; derivações completas em \textcite{kingma2014}.
\end{itemize}

\section{O Modelo FactorVAE}
\label{sec:factorvae}

\begin{itemize}
    \item \textbf{Posição na trajetória:} \textit{FactorVAE} \parencite{duan2022} é a extensão probabilística do Conditional Autoencoder \parencite{gu2021} --- preserva a estrutura fatorial e a não-linearidade aprendida, e acrescenta a representação probabilística dos fatores latentes via inferência variacional.
    \item \textbf{Ideia central:} os retornos de cada ativo são modelados como soma de um componente idiossincrático (gaussiano) e um componente sistemático (combinação linear de fatores latentes); todos os parâmetros da decomposição são saídas de redes neurais que processam exclusivamente informação histórica disponível até $t$.
    \item \textbf{Esquema professor-aluno (visão conceitual):} durante o treino, o \emph{codificador} (professor) acessa os retornos contemporâneos $y_t$ e produz uma posterior informada $q(z|x_t, y_t)$; o \emph{preditor} (aluno) acessa apenas o histórico $x_t$ e produz uma prior $p(z|x_t)$; a função objetivo força o preditor a imitar o professor, transferindo o sinal preditivo. Na inferência fora da amostra, retornos contemporâneos são indisponíveis --- o codificador é descartado; apenas preditor e decodificador participam.
    \item \textbf{Composição gaussiana analítica:} a estrutura linear com hipótese gaussiana permite derivar a distribuição preditiva dos retornos em forma fechada --- média e variância por ativo sem Monte Carlo; a incerteza idiossincrática e a incerteza sobre os fatores se propagam analiticamente.
    \item \textbf{O que não entra aqui:} a formalização matemática completa (arquitetura dos quatro blocos, equações do extrator, codificador, preditor, decodificador, função objetivo) é apresentada no Capítulo~\ref{cap:metodologia}, que torna o modelo operacionalmente replicável.
\end{itemize}

% PLACEHOLDER: figura conceitual da arquitetura --- quatro blocos (Extrator de Características,
% Codificador, Preditor, Decodificador) com indicação visual dos caminhos de treino
% (todos os quatro blocos + retornos contemporâneos $y_t$) vs.\ inferência (apenas
% Extrator, Preditor, Decodificador). Adaptada de \textcite{duan2022}.
